% !TeX program = xelatex
\documentclass{article}
\usepackage[margin=1in]{geometry}
\usepackage{fontspec}
\setmainfont{Arial}
\usepackage{amsmath}
\usepackage{booktabs}
\begin{document}

\title{MTE 380 Milestone 3 Memo - Design Calculations\\Control System (Architecture + Fuzzy Scheduling)}
\author{}
\date{}
\maketitle

\section*{Memo Summary}
\noindent\textbf{Purpose of document:} detail the software architecture and control calculation.

\noindent\textbf{Purpose of calculation:} determine how sensor inputs (IR error and camera line visibility) are converted to adaptive speed limits and PID gains, then convert into bounded wheel commands for motor control.

\noindent\textbf{Result of calculation:} the controller computes normalized inputs $X_1$ and $X_2$, evaluates fuzzy rules, and outputs mode-dependent $(v_{\max}, K_p, K_i, K_d)$ used in motor commands.

\noindent\textbf{Final control output:} left and right motor-drive outputs are computed from the PID-parameterized control law to drive the motors through the motor driver.

\noindent\textbf{Design-decision support:} this calculation supports the team's decision to use a state machine with adaptive control, instead of one fixed PID tuning for every mission phase.

\section{Objective / Purpose Statement}
\noindent\textbf{Design objective:} maximize line-following speed while maintaining stable tracking and preventing line loss.

\noindent\textbf{Design constraints:} limited camera lookahead, bounded motor commands, and phase-dependent mission tasks (rescue, drop-off, homing).

\noindent\textbf{Design decision:} use a state-machine architecture with fuzzy gain scheduling so each line-following phase can reuse the same control logic with adaptive parameters.

\section{References}
\begin{itemize}
\item G. Eleftheriou et al., ``A Fuzzy Rule-Based Control System for Fast Line-Following Robots,'' \textit{DCOSS 2020}, pp. 388--395, DOI: 10.1109/DCOSS49796.2020.00068.
\item E. H. Mamdani, ``Application of Fuzzy Algorithms for Control of Simple Dynamic Plant,'' \textit{Proc. IEE}, 121(12), 1585--1588, 1974, DOI: 10.1049/piee.1974.0328.
\item Pololu, QTR Reflectance Sensors documentation (readLine weighted-average position), accessed Feb. 2026.
\end{itemize}

\section{Design Inputs and Specifications}
\subsection*{3.1 Design Inputs}
\begin{itemize}
\item $x_{j,k}$: commanded speed of motor $j$ at control step $k$.
\item $|X|$: maximum motor speed capability (rpm or normalized full-scale command).
\item $L_k$: detected visible line length from camera ROI.
\item $|L|$: maximum detectable line length in the selected ROI.
\item $e$: lateral line-position error from IR sensor array.
\item $\Delta t$: controller sampling period (s).
\end{itemize}

\subsection*{3.2 Design Specifications}
\begin{center}
\begin{tabular}{lll}
\toprule
Specification & Symbol / Parameter & Value / Range \\
\midrule
State-machine phases & START to END & 6 mission states + decision node \\
Fuzzy speed input & $X_1$ & 0 to 100 \\
Fuzzy line input & $X_2$ & 0 to 100 \\
Defuzzified output & $X^*$ & 1 to 100 (clamped) \\
Camera ROI (default) & \texttt{roi\_bottom\_ratio} & 0.60 \\
Row occupancy threshold & \texttt{row\_occupancy\_frac} & 0.02 \\
Motor command saturation & $u_L$, $u_R$ & -1 to 1 \\
\bottomrule
\end{tabular}
\end{center}

\section{Assumptions}
\begin{itemize}
\item Camera segmentation in the bottom ROI is reliable under operating lighting.
\item IR array provides a usable centroid error for PID correction in all line-following states.
\item Commanded speed is used as a proxy for actual speed when wheel encoders are unavailable.
\item Initial integral gain is set to $K_i=0$ to reduce wind-up risk; tuning can be updated after testing.
\end{itemize}

\section{Analytical Method and Computations}
\subsection*{5.1 Architecture Method (State Machine)}
Mission flow: \texttt{START -> LINE\_READING -> RESCUE -> GOTO\_ZONE -> RESCUE\_DEST\_ARRIVED -> HOMING -> END}.
\begin{itemize}
\item \texttt{START}: initialize sensors, control gains, and mission flags.
\item \texttt{LINE\_READING}: follow the line using IR error feedback plus camera-based fuzzy gain scheduling.
\item Decision ``line ended?'': if false, remain in \texttt{LINE\_READING}; if true, advance to rescue sequence.
\item \texttt{RESCUE}: stop/slow as needed and enable pickup actuator (vacuum on).
\item \texttt{GOTO\_ZONE}: resume line-following control while carrying payload toward drop-off zone.
\item \texttt{RESCUE\_DEST\_ARRIVED}: disable pickup actuator (vacuum off) and confirm drop-off completion.
\item \texttt{HOMING}: follow return route with the same line controller under end-of-mission constraints.
\item \texttt{END}: stop motors and terminate mission logic safely.
\end{itemize}

\subsection*{5.2 Underlying Engineering Principle for Controls}
The underlying engineering principle is nonlinear closed-loop feedback control with sensor fusion and gain scheduling:
\begin{itemize}
\item \textbf{Closed-loop error regulation:} the IR array estimates lateral position error relative to line center, and PID control continuously reduces that error.
\item \textbf{Gain scheduling:} controller aggressiveness is not fixed. It is adapted by fuzzy logic according to operating condition indicators ($X_1$ speed demand and $X_2$ visible line length).
\item \textbf{Sensor fusion by role separation:} IR sensing provides high-rate local tracking error ($e$), while camera sensing provides lookahead visibility ($L_k$). Combining both improves robustness against either sensor's individual weaknesses.
\item \textbf{Actuator-constrained control:} the computed commands are bounded and mapped to motor-driver voltage/PWM commands, enforcing physical limits and reducing saturation-induced instability.
\end{itemize}

\subsection*{5.3 Control Computation Mathematics}
\noindent\textbf{Step 1: architecture-level control decomposition (from paper):}\par
\[
x(k)=F_1(\hat{y}(k))+F_2(\dot{y}(k))
\]
\noindent where $F_1$ uses IR-array information for steering correction and $F_2$ uses camera lookahead information for speed/gain scheduling.

\noindent\textbf{Step 2: IR-array feedback error (closed-loop variable):}\par
\[
e_k=\left(\frac{\sum_{i=1}^{n} i\,s_i}{\sum_{i=1}^{n} s_i}\right)-c
\]
\noindent where $s_i$ is the IR signal at sensor $i$ and $c$ is the center index. This error is the primary feedback variable regulated by PID.

\noindent\textbf{Step 3: camera line-length calculation and normalization inputs:}\par
\noindent Camera visibility is computed from ROI row occupancy:
\[
L_k=\text{rows with line pixels in ROI},\qquad |L|=\text{ROI height in rows}
\]
\noindent Normalized fuzzy inputs:
\[
X_1=\frac{\max_j x_{j,k}}{|X|}\times100,\qquad
X_2=\frac{L_k}{|L|}\times100
\]

\noindent\textbf{Step 4: fuzzification model:}\par
\noindent Crisp inputs $(X_1,X_2)$ are mapped to graded memberships in $[0,1]$.
\noindent For standard triangular sets ($a<b<c$):
\[
\mu_{\triangle}(x;a,b,c)=
\begin{cases}
0,& x\le a \text{ or } x\ge c\\
\dfrac{x-a}{b-a},& a<x<b\\
\dfrac{c-x}{c-b},& b\le x<c
\end{cases}
\]
\noindent For shoulder sets (membership functions with a flat value of 1 on one side):
\[
\mu_{\text{LS}}(x;a,c)=\max\!\left(0,\min\!\left(1,\frac{c-x}{c-a}\right)\right),\quad
\mu_{\text{RS}}(x;a,c)=\max\!\left(0,\min\!\left(1,\frac{x-a}{c-a}\right)\right)
\]

\noindent Input membership sets used by the scheduler are:
\begin{center}
\begin{tabular}{lll}
\toprule
Variable & Label & Triangle $(a,b,c)$ \\
\midrule
$X_1$ & Low & $(0,0,40)$ \\
$X_1$ & Medium & $(20,50,80)$ \\
$X_1$ & High & $(60,100,100)$ \\
$X_2$ & Close & $(0,0,50)$ \\
$X_2$ & Far & $(30,100,100)$ \\
\bottomrule
\end{tabular}
\end{center}
\noindent Apply shoulders where endpoints are repeated:
\[
\mu_{X_1}^{\text{Low}}=\mu_{\text{LS}}(X_1;0,40),\ 
\mu_{X_1}^{\text{High}}=\mu_{\text{RS}}(X_1;60,100),\ 
\mu_{X_2}^{\text{Close}}=\mu_{\text{LS}}(X_2;0,50),\ 
\mu_{X_2}^{\text{Far}}=\mu_{\text{RS}}(X_2;30,100),
\]
\noindent and use $\mu_{\triangle}$ for interior sets (e.g., Medium).

\noindent Output labels represent control modes, each with its own output membership function:
\[
\text{LC}(0,10,25),\ \text{LF}(15,30,45),\ \text{MC}(35,50,65),\ 
\text{MF}(55,70,85),\ \text{HC}(70,85,95),\ \text{HF}(85,100,100)
\]
\noindent Use shoulder handling for repeated endpoints in output sets; specifically,
\[
\mu_{\text{HF}}(x)=\mu_{\text{RS}}(x;85,100),
\]
\noindent while LC, LF, MC, MF, and HC use $\mu_{\triangle}$.
\noindent These output-label shapes are the targets used by the rule table in Step 5 and combined into one final value in Step 6.
\noindent This step lets one input value belong partly to more than one label, so control changes are gradual instead of abrupt.

\noindent\textbf{Step 5: rule base and firing strengths:}\par
\begin{center}
\begin{tabular}{lll}
\toprule
$X_1$ & $X_2$ & Output \\
\midrule
Low & Close & LC \\
Low & Far & LF \\
Medium & Close & MC \\
Medium & Far & MF \\
High & Close & HC \\
High & Far & HF \\
\bottomrule
\end{tabular}
\end{center}
\noindent The table is applied row by row. If row $r$ has input labels $(A_r,B_r)$ and output label $L_r$, then
\[
\alpha_r=\min\left(\mu_{X_1}^{A_r}(X_1),\mu_{X_2}^{B_r}(X_2)\right)
\]
is the row strength for that exact row.
For example, the row ``Medium \& Far $\rightarrow$ MF'' gives
\[
\alpha_{\text{MF}}=\min\left(\mu_{X_1}^{\text{Medium}}(X_1),\mu_{X_2}^{\text{Far}}(X_2)\right).
\]
If multiple rows point to the same output label, keep the largest row strength for that label:
\[
\mu_{\text{rule}}(L)=\max_{r\in L} \alpha_r
\]

\noindent\textbf{Step 6: aggregate outputs, defuzzify, and choose one label:}\par
\noindent For each output label $L\in\{\mathrm{LC,LF,MC,MF,HC,HF}\}$, let $\beta_L=\mu_{\text{rule}}(L)$.
\[
\tilde{\mu}_L(x)=\min(\beta_L,\mu_L(x)),\qquad
\mu_{\text{agg}}(x)=\max_L \tilde{\mu}_L(x)
\]
\[
X^*=\frac{\sum_{x=0}^{100}\mu_{\text{agg}}(x)\,x}{\sum_{x=0}^{100}\mu_{\text{agg}}(x)},\qquad
X_q^*=\min(100,\max(1,\mathrm{round}(X^*,2)))
\]
\[
L^*=\arg\min_{L\in\{\mathrm{LC,LF,MC,MF,HC,HF}\}}|X_q^*-b_L|
\]
\noindent where $b_L$ is the peak value of label $L$.
\noindent This step converts all fuzzy-rule outputs into one scheduler result and one control-mode label.

\noindent\textbf{Step 7: map selected label to speed cap and PID gains:}\par
\[
(v_{\max},K_p,K_i,K_d)=\Theta(L^*),\qquad
v=\min(v_{\text{cmd}},v_{\max})
\]
\begin{center}
\begin{tabular}{lcccc}
\toprule
Label & $v_{\max}$ & $K_p$ & $K_i$ & $K_d$ \\
\midrule
LC & 0.30 & 0.80 & 0.00 & 0.10 \\
LF & 0.35 & 0.70 & 0.00 & 0.10 \\
MC & 0.45 & 0.65 & 0.00 & 0.12 \\
MF & 0.55 & 0.55 & 0.00 & 0.14 \\
HC & 0.65 & 0.45 & 0.00 & 0.16 \\
HF & 0.75 & 0.40 & 0.00 & 0.18 \\
\bottomrule
\end{tabular}
\end{center}

\noindent\textbf{Step 8: compute PID correction and bounded wheel commands:}\par
\[
I_k=I_{k-1}+e_k\Delta t,\qquad
\dot e_k=\frac{e_k-e_{k-1}}{\Delta t}
\]
\[
u_k=K_p e_k + K_i I_k + K_d\dot e_k
\]
\[
u_{L,\text{raw}}=v-u_k,\qquad
u_{R,\text{raw}}=v+u_k
\]
\[
u_L=\text{clamp}(u_{L,\text{raw}},-1,1),\qquad
u_R=\text{clamp}(u_{R,\text{raw}},-1,1).
\]
\noindent This converts scheduled gains and current error into safe, bounded wheel commands.

\section{Calculations}
\subsection*{6.1 Worked numerical example (one control cycle)}
Assume one control cycle with
\[
x_{L,k}=0.62,\ x_{R,k}=0.68,\ |X|=1.00,\ L_k=78\ \text{rows},\ |L|=120\ \text{rows},
\]
\[
e_{k-1}=0.08,\ e_k=0.12,\ \Delta t=0.02\ \text{s},\ v_{\text{cmd}}=0.70.
\]
Normalized inputs:
\[
X_1=\frac{\max(0.62,0.68)}{1.00}\times100=68,\qquad
X_2=\frac{78}{120}\times100=65.
\]
Non-zero memberships and corresponding non-zero rule strengths (all others are 0):
\[
\mu_{\text{Med}}(X_1)=\frac{80-68}{80-50}=0.40,\quad
\mu_{\text{High}}(X_1)=\frac{68-60}{100-60}=0.20,\quad
\mu_{\text{Far}}(X_2)=\frac{65-30}{100-30}=0.50,
\]
\[
\alpha_{\text{MF}}=\min(0.40,0.50)=0.40,\qquad
\alpha_{\text{HF}}=\min(0.20,0.50)=0.20.
\]
\[
X^*=75.30,\qquad X_q^*=75.30,\qquad
L^*=\arg\min_{L}|X_q^*-b_L|=\text{MF}.
\]
Select MF parameters from Step 7:
\[
v_{\max}=0.55,\qquad K_p=0.55,\qquad K_i=0,\qquad K_d=0.14.
\]
PID output and wheel commands:
\[
u_k=0.55(0.12)+0+0.14\frac{0.12-0.08}{0.02}=0.346,
\]
\[
v=\min(0.70,0.55)=0.55,\qquad
u_L=\text{clamp}(0.55-0.346,-1,1)=0.204,
\]
\[
u_R=\text{clamp}(0.55+0.346,-1,1)=0.896.
\]
\section{Results / Conclusions}
The full calculation chain now provides a complete numeric path from sensed inputs to actuator commands: in the worked cycle, the scheduler selects MF, caps speed at $v_{\max}=0.55$, and outputs $u_L=0.204$ and $u_R=0.896$ (both within the $[-1,1]$ actuator bounds). This directly supports the objective of fast but stable line following and supports the design decision to use state-machine mission control with fuzzy gain scheduling, rather than one fixed PID tuning across all mission phases.

\end{document}
